\chapter{Discusión y conclusiones}
\label{cap:discusion}
%% ============================================================================
%% BORRADOR - Capítulo 5
%% Los comentarios que empiezan con "%% CAMI" son notas para revisar antes de
%% la versión final (datos por verificar, figuras por agregar, etc.).
%% ============================================================================

En este capítulo se interpretan los resultados presentados en el Capítulo 4 a la luz de los
objetivos planteados en el Capítulo 2. Primero se discuten los hallazgos sobre la calidad de los
datos, la diversidad alfa y beta, la composición taxonómica y las pruebas estadísticas; después se
señalan las limitaciones del estudio y se presentan las conclusiones. Las perspectivas y el trabajo
a futuro que se derivan de estas conclusiones se presentan en el Capítulo \ref{cap:perspectivas}.

%% ----------------------------------------------------------------------------
\section{Discusión}
%% ----------------------------------------------------------------------------

\subsection{La calidad de los datos condiciona la interpretación}

El primer resultado del análisis fue metodológico: la diversidad alfa calculada sobre los datos
crudos no permitía distinguir entre plantas saludables y no saludables, mientras que, tras eliminar
las cinco muestras con menos de 25 millones de lecturas (MP2079, MP2080, MP2088, MP2109 y MP2137),
las tendencias entre grupos se hicieron visibles. Casos como el de la muestra MP2088, con solo dos
lecturas, muestran que una única muestra defectuosa puede dominar la escala de los índices y ocultar
la estructura del resto de los datos. El filtrado de calidad no es, por lo tanto, un paso
accesorio, sino una condición necesaria para que las comparaciones posteriores tengan sentido.

Sin embargo, aun después del filtrado persistió una diferencia sistemática en la profundidad de
secuenciación: las muestras de plantas saludables tuvieron, en promedio, más lecturas que las de
plantas no saludables. Los índices que dependen del número de taxones observados o de los
\emph{singletons} y \emph{doubletons}, como Chao1, son sensibles a esta diferencia, porque una
muestra con más lecturas tiene mayor probabilidad de detectar taxones raros. Por ello, parte de la
mayor diversidad observada en las plantas saludables podría explicarse por el esfuerzo de muestreo
y no únicamente por el estado de salud de la planta \citep{mcmurdie2014waste}. La rarefacción a
profundidad común permitió separar los dos efectos (Sección \ref{sec:res-normalizacion}): las curvas de rarefacción
muestran que todas las muestras están saturadas, y al igualar el esfuerzo de muestreo la diferencia
de riqueza entre grupos se redujo a menos de la mitad, mientras que la de Shannon, que no depende
de la profundidad, se mantuvo idéntica. Es decir, la mayor riqueza aparente de las plantas saludables
era en buena parte un efecto de la profundidad; la tendencia en equidad, aunque débil, es genuina.
El PERMANOVA tampoco cambió de conclusión con ninguna de las tres normalizaciones probadas, lo que
indica que la ausencia de separación entre grupos no es un artefacto de la estrategia de
normalización \citep{weiss2017normalization}.

\subsection{Diversidad alfa: una tendencia consistente pero débil}

En todos los índices de diversidad alfa (Chao1, Shannon y Simpson) las plantas saludables
mostraron valores ligeramente mayores que las no saludables, y esta tendencia se repitió al separar
los datos por reino y al descender a los niveles de filo, familia, género y especie. Esta dirección
coincide con lo reportado en otros estudios sobre la rizósfera de fresa y de otros cultivos, en los
que la enfermedad se asocia con una pérdida de diversidad microbiana
\citep{xiong2020rhizosphere, gao2021disease, yang2020comparison}.

No obstante, la magnitud de la diferencia fue pequeña. Para el índice de Shannon calculado sobre la
comunidad completa, las medias fueron 7.478 en el grupo saludable y 7.417 en el no saludable, una
diferencia de apenas 0.061 unidades. Ni la prueba t con varianzas distintas ($p = 0.1501$) ni la
prueba de Wilcoxon-Mann-Whitney ($p = 0.2586$) permitieron rechazar la hipótesis nula de igualdad
entre grupos. La tendencia visual, por consistente que sea, no alcanzó a separarse del ruido propio
de los datos.

Un resultado que merece atención es la diferencia en la dispersión de los grupos. La desviación
estándar del índice de Shannon fue de 0.061 en las plantas saludables y de 0.168 en las no
saludables, y la prueba F rechazó con claridad la igualdad de varianzas ($F = 0.132$,
$p < 0.001$). Es decir, las plantas no saludables no difieren tanto en su diversidad promedio como
en su variabilidad: sus comunidades microbianas son más heterogéneas entre sí. Este patrón es
consistente con el llamado \emph{principio de Anna Karenina} aplicado a microbiomas, según el cual
los microbiomas de hospederos sanos tienden a parecerse entre sí, mientras que cada hospedero
enfermo o estresado se desestabiliza a su manera \citep{zaneveld2017stress}. El mismo principio se ha
confirmado recientemente en microbiomas vegetales bajo estrés por patógenos, a partir de más de mil
muestras de maíz en las que las plantas enfermas mostraron una dispersión beta sistemáticamente mayor
\citep{li2026anna}, y la mayor heterogeneidad de las plantas no saludables de esta tesis se verificó
también en la composición completa mediante PERMDISP ($p = 0.023$, Capítulo 4). Desde el punto de
vista estadístico, esta observación sugiere que comparar solo medias puede ser insuficiente y que la
dispersión es, en sí misma, una característica diferenciadora entre los grupos.

\subsection{Diversidad beta: gran superposición entre grupos}

Ninguna de las métricas de distancia utilizadas (Bray-Curtis, Jaccard, Euclidiana, Manhattan y
Jensen-Shannon) produjo una separación visible entre muestras saludables y no saludables en las
ordenaciones. Para cuantificar esta observación se aplicó un análisis PERMANOVA sobre la matriz de
disimilitud de Bray-Curtis de las abundancias relativas: el estado de la planta explicó solo el
2.4\% de la variación total en la composición ($R^2 = 0.024$, $F = 1.26$, $p = 0.167$, 999
permutaciones) \citep{anderson2001new, anderson2011navigating}. Es decir, más del 97\% de la
variación entre muestras se debe a factores distintos del estado sano o enfermo de la planta. La
prueba PERMDISP \citep{anderson2006distance}, en cambio, sí fue significativa ($p = 0.023$): las
plantas no saludables están más dispersas alrededor de su centroide. Esta combinación, centroides
iguales y dispersiones distintas, es precisamente la que recomiendan examinar las revisiones
metodológicas antes de interpretar un PERMANOVA con grupos desbalanceados
\citep{kleinebardenhorst2021data, aldirawi2023univariate}.

Esta superposición indica que la composición global del microbioma rizosférico está dominada por
otros factores, probablemente el tipo de suelo, el manejo agronómico, el sitio y el momento de
muestreo, que no fueron registrados en los metadatos disponibles. Al restringir el análisis a
niveles taxonómicos más finos, en particular género y especie, las diferencias en Bray-Curtis
resultaron más evidentes, lo que sugiere que la señal asociada a la salud de la planta, si existe,
no está en la estructura global de la comunidad sino en un subconjunto reducido de taxones.

\subsection{Composición taxonómica: taxones candidatos, no biomarcadores}

A nivel de filo, Actinobacteria y Proteobacteria fueron los grupos más abundantes en ambos tipos de
muestras, lo que concuerda con la composición típica de la rizósfera de fresa
\citep{yang2020comparison, xiong2020rhizosphere}. La exploración por niveles sugería un cuadro
familiar en la literatura: géneros benéficos como \textit{Bacillus} y \textit{Pseudomonas} más
presentes en las plantas saludables, y fitopatógenos como \textit{Fusarium} y \textit{Ralstonia} en
las no saludables, tal como se ha reportado para fresa invadida por \textit{Fusarium}
\citep{yang2023fusarium}, en el metaanálisis de rizósferas con marchitez \citep{su2025general} y en
los microbiomas supresores de enfermedad \citep{spooren2024plant}. La cuantificación de las
proporciones (Sección \ref{sec:res-candidatos}) no confirmó ese cuadro. \textit{Fusarium} representa alrededor del 0.14\%
de las lecturas en ambos grupos, con distribuciones casi idénticas ($p = 0.978$); \textit{Pseudomonas},
\textit{Bacillus}, \textit{Streptomyces} y \textit{Ralstonia} tienen medianas prácticamente iguales entre
grupos; y la única diferencia nominal, la de \textit{Phytophthora} y su filo Oomycota, va en la
dirección contraria a la esperada para un patógeno, es pequeña y no sobrevive a la corrección por
comparaciones múltiples.

Este resultado merece dos comentarios. El primero es metodológico: las impresiones que dejan las
gráficas de barras con aglomeración al 10\% son engañosas, porque reescalan fracciones minúsculas y
porque unas pocas muestras atípicas, como las dos plantas marchitas con el triple de \textit{Fusarium},
dominan la percepción visual. Solo la comparación cuantitativa, con sus pruebas y su corrección por
comparaciones múltiples, permite decidir. El segundo es biológico: que \textit{Fusarium} no esté
enriquecido en la rizósfera de las plantas marchitas no descarta su papel en la enfermedad. La
clasificación con Kraken asigna lecturas al género y no distingue cepas patógenas de saprófitas; la
marchitez puede deberse a otros agentes, como \textit{Verticillium} o \textit{Macrophomina}, no
necesariamente bien representados en la base de datos; y la rizósfera no es el tejido donde el
patógeno se concentra. Además, los datos son composicionales: las abundancias relativas suman uno,
de modo que el aumento de un taxón obliga a la disminución aparente de otros, aun cuando su
abundancia absoluta no haya cambiado \citep{gloor2017microbiome, calle2019statistical}. En estos
datos, por lo tanto, no se identifican géneros biomarcadores del estado de salud; lo que sí se
identifica es que la señal, si existe, no está en la abundancia de un taxón particular.

\subsection{Diversidad de eucariotas a nivel de género: una señal que no se confirmó}

Se esperaba que la señal asociada con los fitopatógenos, diluida en la comunidad completa, resultara
detectable al aislar la fracción eucariota, y en particular al índice de Shannon calculado sobre los
58 géneros de eucariotas, grupo que incluye a \textit{Fusarium} y a \textit{Phytophthora}. No fue
así: las medias fueron 3.077 en las plantas saludables y 3.035 en las no saludables, y ni la prueba
t ($p = 0.260$ con varianzas iguales, $p = 0.378$ con Welch) ni la de Wilcoxon-Mann-Whitney
($p = 0.963$) rechazaron la igualdad entre grupos.

Este resultado tiene una historia que conviene dejar registrada. Una versión anterior del análisis
reportaba para este mismo contraste una diferencia significativa ($p = 0.0017$), que llegó a
presentarse como el hallazgo principal del trabajo. Al recalcularla se encontró que el script
construía la tabla de abundancias con funciones que reordenan las muestras por abundancia y luego la
unía a los metadatos sin emparejar por nombre de muestra, de modo que cada valor de diversidad
quedaba asignado al estado de salud de otra planta. La ``diferencia'' era un artefacto de ese
desemparejamiento. El error es instructivo desde el punto de vista del análisis de datos: una
verificación tan simple como comprobar que los identificadores de muestra coinciden antes de unir
dos tablas habría evitado un resultado espurio con un valor $p$ muy convincente. Por eso todas las
cifras del Capítulo 4 se recalcularon con scripts que verifican explícitamente ese emparejamiento.

Descartada esa señal, la conclusión es coherente con el resto del capítulo: la comunidad eucariota
tampoco distingue a los grupos en su nivel medio de diversidad. Cabe recordar, además, que este
contraste fue uno de varios realizados sobre distintos subconjuntos y niveles taxonómicos, de modo
que cualquier resultado significativo que aparezca en análisis futuros deberá evaluarse con una
corrección por comparaciones múltiples, por ejemplo la de Benjamini-Hochberg
\citep{benjamini1995controlling}.

\subsection{Tamaño de muestra y potencia estadística}

La ausencia de significancia en la comunidad completa no debe leerse como evidencia de que no exista
diferencia entre los grupos. Con las medias y desviaciones estándar observadas del índice de
Shannon, el tamaño del efecto estandarizado (\emph{d} de Cohen) fue de aproximadamente 0.56, un
efecto de magnitud moderada. Con los 35 y 18 metagenomas disponibles, la potencia de una prueba t
para detectar un efecto de ese tamaño con $\alpha = 0.05$ es de apenas 0.48 aproximadamente: aun si
la diferencia fuera real, el diseño la detectaría menos de la mitad de las veces. Para alcanzar una
potencia de 0.8 se necesitarían alrededor de 51 muestras por grupo, casi el triple del número
actual de plantas no saludables \citep{xia2018statistical}.
%% CAMI: potencia calculada con power.t.test usando la media armónica de los tamaños de grupo
%% (n aprox. 24) y la desviación estándar combinada (0.109). Es una aproximación, dado que las
%% varianzas no son iguales; puede incluirse como resultado en el Capítulo 4.

Este cálculo reencuadra los resultados del Capítulo 4: la falta de significancia es compatible con
un efecto real que el diseño no tiene potencia suficiente para detectar. El desbalance entre grupos,
con casi el doble de muestras saludables que no saludables, y la mayor varianza precisamente en el
grupo más pequeño agravan el problema. Esta limitación motiva directamente el trabajo a futuro
descrito en el Capítulo \ref{cap:perspectivas}.

\subsection{Limitaciones del estudio}

Además del tamaño de muestra, el estudio presenta las siguientes limitaciones:

\begin{itemize}
    \item \textbf{Diseño observacional y metadatos limitados.} Solo se contó con la etiqueta
    saludable o no saludable de cada planta. No se dispuso de variables como tipo de suelo, sitio,
    fecha de muestreo, variedad o manejo agronómico, que podrían explicar buena parte de la
    variación observada en la diversidad beta.
    \item \textbf{Etiquetado binario de la enfermedad.} La categoría ``no saludable'' agrupa
    plantas marchitas sin identificar el agente causal ni la severidad de la enfermedad, lo que
    puede contribuir a la alta variabilidad observada en ese grupo.
    \item \textbf{Clasificación taxonómica.} La asignación con Kraken depende de la base de datos
    de referencia; los eucariotas, en particular los hongos, están menos representados que las
    bacterias, lo que puede sesgar su detección y abundancia.
    \item \textbf{Profundidad de secuenciación desigual.} Las muestras saludables tuvieron más
    lecturas, lo que puede inflar los índices de riqueza en ese grupo.
    \item \textbf{Asociación y no causalidad.} Las diferencias observadas son asociaciones; con el
    diseño actual no es posible determinar si los cambios en el microbioma son causa o consecuencia
    de la enfermedad.
\end{itemize}

%% ----------------------------------------------------------------------------
\section{Conclusiones}
%% ----------------------------------------------------------------------------

En relación con los objetivos planteados, se concluye lo siguiente:

\begin{enumerate}
    \item \textbf{Caracterización de la composición y la diversidad.} Se caracterizó el microbioma
    rizosférico de 53 plantas de fresa, 35 saludables y 18 no saludables, a los niveles de reino,
    filo, familia, género y especie. La comunidad está dominada por bacterias, principalmente
    Actinobacteria y Proteobacteria, en ambos grupos. Las plantas saludables mostraron de forma
    consistente una diversidad alfa ligeramente mayor en todos los índices y niveles taxonómicos.
    Además, el filtrado de calidad resultó indispensable para obtener comparaciones interpretables,
    y la rarefacción mostró que la diferencia de riqueza se debe en buena parte a la mayor
    profundidad de secuenciación de las plantas saludables, mientras que la diferencia en Shannon es
    independiente de la profundidad.

    \item \textbf{Evaluación estadística de las diferencias.} Ni en la comunidad completa ni en la
    fracción eucariota las diferencias en diversidad alfa fueron estadísticamente significativas, y
    el estado de la planta explicó una fracción muy pequeña de la variación en la composición. Las
    únicas pruebas que rechazaron su hipótesis nula fueron las de igualdad de dispersión, la F
    sobre las varianzas de Shannon y PERMDISP sobre la composición: las plantas no saludables son
    significativamente más heterogéneas entre sí. El análisis de potencia muestra que
    el tamaño de muestra actual es insuficiente para detectar efectos moderados, por lo que los
    resultados no significativos deben interpretarse como no concluyentes y no como ausencia de
    efecto.

    \item \textbf{Patrones asociados con el estado de salud.} Las plantas no saludables presentaron
    comunidades más heterogéneas entre sí, con una varianza de la diversidad significativamente
    mayor. En cambio, la comparación cuantitativa de las proporciones de \textit{Fusarium} y de
    los géneros candidatos que la exploración había sugerido, \textit{Pseudomonas}, \textit{Bacillus},
    \textit{Streptomyces}, \textit{Paenibacillus}, \textit{Ralstonia} y \textit{Phytophthora}, no mostró
    diferencias entre grupos: ningún género está enriquecido en una condición una vez que se
    corrige por comparaciones múltiples. El estado de salud no se refleja en la abundancia de un
    taxón particular sino en la estabilidad del conjunto, lo que orienta hacia modelos que capturen
    las relaciones entre microorganismos.

    \item \textbf{Clase práctica.} El análisis se convirtió en un taller de 50 minutos, en forma de
    cuaderno de Google Colab de acceso abierto, que reproduce con los mismos datos cada prueba de
    esta tesis y la explica con ejercicios resueltos. El taller enseña a elegir la prueba según el
    diseño y los supuestos, a construir un valor $p$ por permutaciones y a reportar tamaño del
    efecto, estadístico y valor $p$, y recoge como lecciones los dos problemas que más pesaron en
    este trabajo: el emparejamiento de tablas por identificador y la dependencia del valor $p$
    respecto de los supuestos de cada prueba.
\end{enumerate}

En conjunto, este trabajo muestra que el análisis de datos metagenómicos requiere algo más que la
aplicación directa de índices ecológicos: exige controlar la calidad y la profundidad de los datos,
tener en cuenta su naturaleza composicional, evaluar la potencia estadística del diseño y considerar
la dispersión, y no solo las medias, como una característica informativa. Las tendencias observadas
son consistentes con la hipótesis de que la salud de la planta de fresa se asocia con un microbioma
rizosférico más diverso y estable, pero confirmarla requiere más datos y modelos capaces de capturar
la estructura de dependencia entre los microorganismos.
